\section{The Conformal Group in Two Dimensions}

If the dimension of spacetime is larger than two (\(D>2\)), the conformal group is a finite-dimensional Lie group. While it is larger than the standard Poincar\'e group and therefore provides more constraints on the theory, the number of these constraints remains strictly finite. However, in exactly two dimensions, the situation changes drastically and the conformal symmetry becomes infinite-dimensional.

To see this, let us work in 2D Euclidean space (if we start from a Minkowski metric, we can always perform a Wick rotation). In a 2D Minkowski spacetime, the natural variables would be the \textit{lightcone coordinates}; in Euclidean space, the analogous operation is to map the 2D Cartesian plane into the complex plane.
We take a new basis by changing coordinates from the real Cartesian coordinates \(x^1\) and \(x^2\) to the complex coordinates \(z\) and \(\overline{z}\):
\[
    \begin{dcases}
        z = x^1 + i x^2, \\
        \overline{z} = x^1 - i x^2,
    \end{dcases} \implies \begin{dcases}
        x^1 = \frac{z + \overline{z}}{2}, \\
        x^2 = \frac{z - \overline{z}}{2i}.
    \end{dcases}
\]
While physically one coordinate is the complex conjugate of the other, it is mathematically very convenient to treat \(z\) and \(\overline{z}\) as two formally independent variables characterizing a point in \(\mathbb{C}^2\).

In these new coordinates, the flat Euclidean line element becomes:
\[
    \mathrm{d} s^2 = (\mathrm{d} x^1)^2 + (\mathrm{d} x^2)^2 = \mathrm{d} z \mathrm{d} \overline{z},
\]
where \(\mathrm{d} z \mathrm{d} \overline{z}\) also acts as the area element (\(\mathrm{d}^2 x\)) up to a constant factor. The corresponding derivatives with respect to the complex coordinates are:
\[
    \begin{dcases}
        \partial_z = \frac{1}{2}(\partial_1 - i \partial_2), \\
        \partial_{\overline{z}} = \frac{1}{2}(\partial_1 + i \partial_2),
    \end{dcases} \implies \partial_z z = 1, \quad \partial_z \overline{z} = 0,
\]
and similarly for \(\partial_{\overline{z}}\). From the line element \(\mathrm{d}s^2 = g_{\mu\nu}\mathrm{d}x^\mu \mathrm{d}x^\nu\), we can read off the metric tensor in the \((z, \overline{z})\) basis, which takes the off-diagonal form:
\[
    g = \begin{pmatrix}
        0     & 1 / 2 \\
        1 / 2 & 0
    \end{pmatrix}.
\]
Consequently, any field in our theory can be redefined as a function of these two independent variables:
\[
    \varphi(x) \to \varphi(z, \, \overline{z}).
\]

Now, let us consider conformal transformations. A coordinate transformation is conformal if it simply rescales the line element by a strictly positive local factor \(\Omega\):
\[
    (\mathrm{d}s^{\prime})^2 =  \Omega(z, \, \overline{z}) \mathrm{d} s^2.
\]
In complex coordinates, we can achieve this by applying independent transformations to \(z\) and \(\overline{z}\). Let us map them to new variables \(w\) and \(\overline{w}\):
\[
    \begin{dcases}
        z \mapsto w = f(z), \\
        \overline{z} \mapsto \overline{w} = \overline{f}(\overline{z}),
    \end{dcases} \quad \implies \quad \begin{dcases}
        \mathrm{d} w = f^{\prime}(z) \mathrm{d} z, \\
        \mathrm{d} \overline{w} = \overline{f}^{\prime}(\overline{z}) \mathrm{d} \overline{z}.
    \end{dcases}
\]
Here, \(f(z)\) is an analytic (holomorphic) function, meaning its derivative \(f'(z)\) exists, and \(\overline{f}(\overline{z})\) is an anti-holomorphic function. The new line element becomes:
\[
    (\mathrm{d} s^{\prime} )^2 = \mathrm{d} w \mathrm{d} \overline{w} = f^{\prime}(z) \overline{f}^{\prime}(\overline{z}) \mathrm{d} z \mathrm{d} \overline{z},
\]
which perfectly matches the conformal condition with the scale factor completely factorized into a holomorphic and an anti-holomorphic part:
\[
    \Omega(z, \, \overline{z}) = f^{\prime}(z) \overline{f}^{\prime}(\overline{z}).
\]
This reveals a profound result: \textit{any} analytic transformation of the complex plane is a conformal transformation. Geometrically, analytic functions preserve local angles, maintaining the shape of infinitesimally small figures, even though globally they can stretch and severely distort areas (a classic example is the Mercator projection of a sphere, where the North Pole is stretched into an entire infinite line).

Since there are infinitely many analytic functions, the conformal group in two dimensions is \textbf{infinite-dimensional}. This implies we have infinitely many symmetries and, consequently, infinitely many constraints on the algebra of local fields. When a theory is subjected to infinitely many constraints, it often becomes exactly solvable, allowing us to compute quantities far beyond standard perturbation theory.

To rigorously link this back to the conformal Killing equation, consider an infinitesimal conformal transformation:
\[
    \begin{aligned}
        z            & \mapsto z + \epsilon(z),                                  \\
        \overline{z} & \mapsto \overline{z} + \overline{\epsilon}(\overline{z}),
    \end{aligned}
\]
where we define the displacement vector components as:
\[
    \epsilon(z) = \epsilon_1 + i \epsilon_2, \quad \overline{\epsilon}(\overline{z}) = \epsilon_1 - i \epsilon_2.
\]
By definition, an infinitesimal transformation is conformal if it satisfies the conformal Killing equation. In \(D=2\) dimensions, this equation reads:
\[
    \partial_\mu \epsilon_\nu + \partial_\nu \epsilon_\mu = \frac{2}{2} (\partial \cdot \epsilon) \delta_{\mu\nu} = (\partial_1\epsilon_1 + \partial_2\epsilon_2) \delta_{\mu\nu}.
\]
Evaluating this tensor equation component by component (\(\mu=1, \nu=1\) and \(\mu=1, \nu=2\)), we obtain exactly the celebrated \textbf{Cauchy-Riemann equations}:
\[
    \begin{dcases}
        \partial_1 \epsilon_1 = \partial_2 \epsilon_2, \\
        \partial_1 \epsilon_2 = -\partial_2 \epsilon_1.
    \end{dcases}
\]
These equations are the necessary and sufficient condition for \(\epsilon(x^1, x^2) = \epsilon_1 + i \epsilon_2\) to be a complex analytic function \(\epsilon(z)\). This confirms that in two dimensions, the set of all conformal transformations coincides precisely with the set of all holomorphic and anti-holomorphic mappings.

    [...]

\[
    \epsilon(z) = \sum_{n = -\infty}^{\infty} \epsilon_n z^{n + 1},
\]
unless you have a constant function, if it is analytic there must be somewhere in the domain of analiticity in which the function is greater or higher than some value, and outside the domain of analiticity the function can diverge. This laurent series is the most general form of an analytic function, and it is convergent in the annulus defined by the inner and outer radius of convergence.

We can consider \(\epsilon_n(z)\) our basis for the expansion
\[
    \epsilon_n(z) = - \epsilon z^{n + 1},
\]
then the generators of the transformations on scalar fields \(\varphi(x)\) as
\[
    \hat{\ell}_n = - z^{n + 1} \partial_z,
\]
such that
\begin{enumerate}
    \item \(\hat{\ell}_0 = -z \partial_z\) is the generator of \textbf{dilatations},
    \item \(\hat{\ell}_{-1} = -\partial_z\) is the generator of \textbf{translations},
    \item \(\hat{\ell}_1 = -z^2 \partial_z\) is the generator of \textbf{special conformal transformations}.
\end{enumerate}

The dilatation or translation is in \(z\), but we need the \(\overline{z}\) direction, such that we need to define
\[
    \overline{\ell}_n = - \overline{z}^{n + 1} \partial_{\overline{z}}.
\]
So we have
\begin{enumerate}
    \item \(\overline{\ell}_0 = -\overline{z} \partial_{\overline{z}}\) is the generator of \textbf{dilatations},
    \item \(\overline{\ell}_{-1} = -\partial_{\overline{z}}\) is the generator of \textbf{translations},
    \item \(\overline{\ell}_1 = -\overline{z}^2 \partial_{\overline{z}}\) is the generator of \textbf{special conformal transformations}.
\end{enumerate}

Thus if we were to write the four momentum operator in the complex coordinates, we would have
\[
    P = \begin{pmatrix}
        E \\
        p
    \end{pmatrix} = \begin{pmatrix}
        \ell_{-1} + \overline{\ell}_{-1} \\
        i(\ell_{-1} - \overline{\ell}_{-1})
    \end{pmatrix},
\]
where \(E\) is the energy and \(p\) is the momentum. The dilatation operator would be
\[
    D = \hat{\ell}_0 + \overline{\ell}_0,
\]

[...]

If you reason in higher dimensions, the special conformal group is six-dimensional, since
\[
    D=2 \implies \frac{(D+1)(D+2)}{2} = 6.
\]

Imagine to apply the commutator \(\left[ \ell_n, \, \ell_m \right]\) on a test function \(f(z)\), we can find
\[
    \left[ \ell_n, \, \ell_m \right] = \left( (n - m) z^{n + m + 1} \partial_z \right) = (n - m) \ell_{n + m},
\]
which is the Witt algebra, the algebra of the generators of conformal transformations in two dimensions. The same holds for the barred generators \(\overline{\ell}_n\). The Witt algebra is an infinite-dimensional Lie algebra, since \(n\) and \(m\) can take any integer value and be applied on any test function \(f(z)\) in the domain of analyticity. The Witt algebra is the classical version of the Virasoro algebra, which is the quantum version of the conformal algebra in two dimensions. Visaroro was a young phd in Turin when he discovered the algebra, after is supervisor (which already had found the Mobius transformations we are going to introduce later) told him to find the algebra of the generators of conformal transformations in two dimensions, and he found this infinite-dimensional algebra.

This generators, with \(n = -1,0,1\) (barred and not), are six generators which form the algebra of the \(2\times 2\) special linear group \(\mathfrak{sl}(2, \, \mathbb{C})\), which is the double cover of the Lorentz group \(SO(3,1)\). This is the conformal group in two dimensions. It is isomorphic to \(SU(2)\oplus SU(2)\).

We have an infinite number of conformal transformations, since they are all the analytic functions, so how can it be six dimensional? We can map the coordintes of the plane to \(z=e^{w}\), which inverse gives the logarithm of \(z\): we are mapping the plane to a strip (given by the riemann surface of the logarithm), and with periodic boundary conditions we get a cylinder. The conformal transformations on the plane are mapped to conformal transformations on the cylinder, and the six generators of the conformal group are mapped to the six generators of the conformal group on the cylinder. very important (it is a way of compactifying the plane).

Now we can ask which are the transformations starting from a plane and ending in a plane, i.e. which are the transformations that do not change the point at infinity. These are the Mobius transformations, which are given by
\[
    w = \frac{az + b}{cz + d}, \quad ad - bc = 1,
\]
which has six parameters. Thus global conformal transformations are given by the Mobius transformations, which are a subgroup of the infinite-dimensional group of local conformal transformations.

    [...]

We can now infer someting about the fields: after the remapping
\[
    \begin{dcases}
        z \\
        \overline{z}
    \end{dcases} \mapsto \begin{dcases}
        w = f(z) \\
        \overline{w} = \overline{f}(\overline{z})
    \end{dcases},
\]
we can find fields that transforms as
\[
    \dots
\]
which are called \textbf{primary fields}, and they are the building blocks of the theory. We can also understand quasi-primary fields, which are fields that transforms as primary fields under the global conformal transformations (Mobius transformations), but not under the local conformal transformations. There are also other fields which do not have tensorial stucture under the conformal transformations, but they are not important for us. Some of the secondary fields happens to be quasi-primary, but not all of them. The primary fields are the most important fields in the theory and we will distiguish them from the secondary fields, which are obtained by applying the generators of the conformal transformations on the primary fields.

For example, the energy-momentum tensor \(T(z)\) is secondary field, and it would seem lika a bad news since it is the most important field in the theory, but we will see the consequences.

If we compute tge derivatives of \(w\) and \(\overline{w}\) with respect to \(z\) and \(\overline{z}\), we can find
\[
    \dots
\]
and now expanding
\[
    \phi(z, \, \overline{z}) = \phi(w - \epsilon, \, \overline{w} - \overline{\epsilon}) = \phi(w, \, \overline{w}) - \epsilon \partial_w \phi(w, \, \overline{w}) - \overline{\epsilon} \partial_{\overline{w}} \phi(w, \, \overline{w}) + \dots
\]
such that we can find the variation of the field \(\phi\) under the conformal transformation
\[
    \delta \phi = - \epsilon \partial_z \phi - \overline{\epsilon} \partial_{\overline{z}} \phi - h \partial_z \epsilon \phi - \overline{h} \partial_{\overline{z}} \overline{\epsilon} \phi,
\]
where we can see how the two variables are completely decoupled, and the transformation of the field is given by the sum of the transformation in \(z\) and the transformation in \(\overline{z}\).

dilatations\dots
\[
    D = \ell_0 + \overline{\ell}_0 = h + \overline{h} = \Delta,
\]

rotations\dots
\[
    L = \ell_0 - \overline{\ell}_0 = h - \overline{h} = s,
\]

such that we can characterize the fields by their scaling dimension \(\Delta\) and their spin \(s\) (or equivalently by \(h\) and \(\overline{h}\)). A scalar field is spinless, so \(s=0\) and \(h = \overline{h}\). A chiral field is a field that depends only on \(z\) or only on \(\overline{z}\), so it has either \(h=0\) or \(\overline{h}=0\).




Thus with any analytic map
\[
    \dots
\]
we can write a general conformal transformation of the field \(\Phi\) as
\[
    \Phi(z, \overline{z}) = \left(\frac{\mathrm{d} w}{\mathrm{d} z} \right)^{-h} \left(\frac{\mathrm{d} \overline{w}}{\mathrm{d} \overline{z}} \right)^{-\overline{h}} \Phi(w, \overline{w}) + \dots
\]
where the first term is the transformation of a primary field (the subset which transforms tensorially), and the dots represent the transformation of the secondary fields. The primary fields are the building blocks of the theory, and all the other fields can be obtained by applying the generators of the conformal transformations on the primary fields.

The new fields \(\Phi(w, \overline{w})\) can be labelled by \((h, \overline{h})\), the power of the Jacobian of the transformation.

\begin{example}[Exercise: transformation of the derivtive of a field]
    \[
        \partial_z \Phi(z, \overline{z}) \mapsto \partial_w \Phi(w, \overline{w})
    \]
    [...]
    we will get some terms which are not the transformation of a primary field (i.e. some second order derivatives of the field), so the derivative of a primary field is not a primary field, in general, but a secondary field.
\end{example}

When we define teh fields on a plane, they are \textbf{quasi-primary fields} if they transform as primary fields under the global conformal transformations (the subgroup of conformal trasnformations which are global, the Moebius transformations), but not under the local conformal transformations: if we try a more general transformation for the energy momentum tensor for example, then it acquires some other terms in the transformation. When we will arrive to the classification of fields in a theory based on the local algebra fields of the thepry (which we can specify), then we will see that there exists some family of fields which are quasi-primary but not primary, and they are called the \textbf{descendants} of the primary fields. The descendants are obtained by applying the generators of the conformal transformations on the primary fields, and they form a representation of the conformal algebra. The primary fields are the highest weight states of this representation, and the descendants are obtained by applying the lowering operators (the generators with \(n > 0\)) on the primary fields.\TODO{check}

The numbers \((h,\overline{h})\) do not label only the primary fields: even secondary or quasi-primary fields can be labeled by \((h,\overline{h})\), since the transformations start always witrh a tensorial components (for the derivative it would be \((h+1, \overline{h})\) for example), which thus express the local scaling properties of the field, even if the field is not primary. The primary fields are the only fields which transform tensorially under the local conformal transformations, but all the fields have well defined scaling properties, which are encoded in the numbers \((h,\overline{h})\).

\[
    z \to  \lambda z, \quad \overline{z} \to \lambda \overline{z} \implies \Phi(z, \overline{z}) \to \lambda^{-h} \lambda^{-\overline{h}} \Phi(z, \overline{z}) = \lambda^{-\Delta} \Phi(z, \overline{z}),
\]
which implies that the scale dimension \(\Delta\) is always given by the sum of the two numbers \(h\) and \(\overline{h}\), and their difference gives the spin \(s\):
\[
    \Delta = h + \overline{h}, \quad s = h - \overline{h}.
\]
Indeed we can think about the generators to notice that
\[
    \Delta = z \partial_z + \overline{z} \partial_{\overline{z}} = h + \overline{h}, \quad s = z \partial_z - \overline{z} \partial_{\overline{z}} = h - \overline{h}.
\]

Thus, we have seen how useful is the choice of complex coordinates in two dimensions, since it allows us to decouple the conformal transformations into two independent sectors, the holomorphic and the anti-holomorphic sector, which are completely decoupled and can be treated separately. But from now on we can adopt the classification in in \(h\) and \(\overline{h}\) (\(\Delta\) and \(s\) practically, the holomorphic and anti-holomorphic scaling dimensions), since they are equivalent.

For scalar fields, we have \(s=0\) and thus \(h = \overline{h}\), and we can expect locally that this number is integer or half-integer. For \textbf{chiral fields}, we have either \(h=0\) or \(\overline{h}=0\), which means that they depend only on one of the two coordinates:
\[
    \begin{dcases}
        \Phi_{h,0}(z) \mapsto \partial_{\overline{z}} \Phi_{h,0}(z) = 0, \quad (s=h) \\
        \Phi_{0,\overline{h}}(\overline{z}) \mapsto \partial_z \Phi_{0,\overline{h}}(\overline{z}) = 0, \quad (s=-h).
    \end{dcases}
\]
Thus we have got basically two continuity equations, one for the holomorphic sector and one for the anti-holomorphic sector, which are completely decoupled. We have to think about some sort of currents, where the spin/helicity gives us a sort of direction...

    [...]

\subsubsection{Stress-Energy Tensor}
We know that the stress-energy tensor \(T_{\mu\nu}\) is a local field, symmetric (or at least it can be made symmetric), and it is a conserved current, so it must satisfy the following conditions:
\[
    T_{\mu\nu}(x) = T_{\nu\mu}(x), \quad \partial^\mu T_{\mu\nu}(x) = 0, \quad T^\mu_\mu(x) = 0,
\]
where the continuity equation \(\partial^\mu T_{\mu\nu} = 0\) expresses the \textit{conservation of energy and momentum}, and the traceless condition \(T^\mu_\mu = 0\) is a consequence of \textit{scale invariance} at criticality (thus for Polyakov, since we have both rototranslation and scale invariance we know it to be conformal invariant).

We can again change to complex coordinates
\[
    \begin{dcases}
        z = x^1 + ix^2, \\
        \overline{z} = x^1 - ix^2,
    \end{dcases}
\]
and we can define the following components of the stress-energy tensor:
\[
    \begin{dcases}
        T(z,\overline{z}) = 2 \pi T_{zz}(z,\overline{z}) = 2 \pi( T_{11} - T_{22} + 2i T_{12} ), \\
        \overline{T}(z,\overline{z}) = 2 \pi T_{\overline{z}\overline{z}}(z,\overline{z}) = 2 \pi( T_{11} - T_{22} - 2i T_{12} ),
        \Theta(z,\overline{z}) = 2 \pi T_{z\overline{z}}(z,\overline{z}) = 2 \pi( T_{11} + T_{22} )=0.
    \end{dcases}
\]
where the \(2\pi\) factor comes from the original work of Polyakov, and it is just a convention. The last component \(\Theta\) is zero due to the traceless condition. The first two components \(T\) and \(\overline{T}\) are called the holomorphic and anti-holomorphic components of the stress-energy tensor, and they are the most important fields in the theory, since they generate the conformal transformations.

We can write the conservation equation \(\partial^\mu T_{\mu\nu} = 0\) in complex coordinates, and we find that it decouples into two independent equations:
\[
    \begin{dcases}
        \partial_{\overline{z}} T_{z z} + \partial_z T_{z\overline{z}} = 0, \\
        \partial_z \overline{T}(\overline{z}) + \dots = 0,
    \end{dcases}
\]
but since \(\Theta\) is zero, we can devolop
\[
    \dots
\]
and find out that \(T\) is undependent of \(\overline{z}\) and \(\overline{T}\) is independent of \(z\):
\[
    T=T(z), \quad \overline{T} = \overline{T}(\overline{z}), \quad \begin{dcases}
        \partial_{\overline{z}} T(z) = 0, \\
        \partial_z \overline{T}(\overline{z}) = 0.
    \end{dcases}
\]

Next let us take into account the action  as the integral of the lagrangian density, which is a local field, and thus we can write
\[
    S = \int \mathrm{d}^2 x \, \mathcal{L}(x) = \int \mathrm{d} z \mathrm{d} \overline{z} \, \mathcal{L}(z, \overline{z}),
\]
and since the action is supposed to have zero mass dimension, the lagrangian density must have scaling dimension \(\Delta = 2\). Thus, the lagrangian density is a local field with scaling dimension 2, and since we derive the stress-energy tensor from the lagrangian through
\[
    T = \frac{\partial \mathcal{L}}{\dots }\dots
\]
then we knwo the stress-energy tensor must have scaling dimension 2 as well, and thus it is a local field with scaling dimension 2. We would have expected it to have three components, but the traceless condition reduces it to two independent components, which are the holomorphic and anti-holomorphic components \(T(z)\) and \(\overline{T}(\overline{z})\).

We know that \(T\) is \((2,0)\) while \(\overline{T}\) is \((0,2)\), so they are not primary fields, since they do not transform tensorially under the conformal transformations, but they are secondary fields...

We know that a two point fucntion at the critical point should rescale by a power law, and thus we can write
\[
    \langle T(z)T(w)\rangle \sim \frac{1}{\vert z- w \vert^4 },
\]
thus if we do a mobius transformation when approaching the 4 oredr pole, it is fine, but a slightly more general transformation will give us some extra terms, which are not the transformation of a primary field, so \(T\) is not a primary field, but it is a secondary field. The same holds for \(\overline{T}\). \uline{Thus the stress-energy tensor is not a primary field.}

Now we are going to show two simple but fundamental results in two dimensional conformal field theory:
\begin{enumerate}
    \item The free boson in two dimensions as a conformal field theory.
    \item The free fermion in two dimensions as a conformal field theory.
\end{enumerate}

\subsection{The Free Boson in Two Dimensions}

[...]
\[
    S = \frac{1}{8\pi} \int \mathrm{d}^2 x \, \partial_\mu \varphi(x) \partial^\mu \varphi(x),
\]
where the factor \(1/8\pi\) is just a convention, choosen by Ginsparg, but it is not important, the calculation are simpler.

    [comment on zero mass term]

A first result we can express with just field theory is
\[
    \langle \varphi(x)\varphi(y) \rangle = - \log \vert x - y \vert^2 + \text{const.},
\]
where \(x,y\) are two dimensional vectors, and the constant is an arbitrary constant which can be fixed by some normalization condition. By changing to complex coordinates, we can write
\[
    \langle \varphi(x)\varphi(y)\rangle = - \log (z - w ) + \log (\overline{z} - \overline{w}) + \text{const.},
\]
effectively decoupling the holomorphic and anti-holomorphic sectors. Furthermore if we consider the d'Alembert equation for the free boson
\[
    \box \varphi(z,\overline{z}) = \dots
\]
so that the more general solution of the equation of motion is given by the sum of a holomorphic and an anti-holomorphic, decoupled.

\begin{remark}
    The presence of the logarithm in the two-point function is a consequence of the fact that the free boson in two dimensions is a massless field, and thus it has a long-range correlation. This is not a conformal field, since we cannot put it in the set of fields of the algebra of local fields of the theory; actually Coleman in the Seventie wrote a paper showing that the massless boson in two dimensions is not a conformal field, but it does not mean that it can be a generator of a bigger algebra of fields, and thus it can be used to construct a conformal field theory, as we will see in the next chapters. If we insist with this lagrangian in thinking that \(\varphi\) as a foundamental field, we are simply wrong, but we can use this to build composite fields, which are conformal fields, and thus we can use the free boson to build a conformal field theory.
\end{remark}

Since the two components are totally decoupled, we can work with \(z\) and all the results will hold for \(\overline{z}\) as well.

We can define the following
\[
    \begin{aligned}
        \langle \partial \phi(z) \partial \phi(w) \rangle = \partial_z \partial_w \langle \phi(z) \phi(w) \rangle \\
        = - \partial_z \partial_w \log (z - w) = - \partial_w \frac{1}{z - w} = - \frac{1}{(z - w)^2},
    \end{aligned}
\]
and since there is a minus sign we can ``be russians about it'' and write the action with a minus sign in front, or instead define the following field
\[
    J(z) = i \partial \phi(z),
\]
for which we can write
\[
    \langle J(z) J(w) \rangle = \frac{1}{(z - w)^2},
\]
which has good properties, and it is a conformal field with scaling dimension \(\Delta = 1\) (since it transforms as the derivative of a primary field). It is a field that depends onaly on \(z\) and it is a \((1,0)\) field, so it is a chiral field. Its spin is \(s=1\) while for \(\overline{J}(\overline{z})\) we would find the same scaling dimension but spin \(s=-1\) and
\(\partial_\mu J^{\mu} = \partial J = \overline{\partial} \overline{J} =\overline{\partial} J = 0\), such that \(J\) and \(\overline{J}\) are conserved currents, and thus they are chiral currents. The charge implemented by the current do not necessary have a physical meaning, but we know that there have to be some symmetries in the theory.

Now let us try to study the operators product expansion of the field \(J(z)\) with itself, which is a fundamental tool in conformal field theory. We can write
\[
    \begin{aligned}
        J(z)J(w) = \contraction{}{J(z)}{}{J(w)} J(z) J(w) + \dots
    \end{aligned}
\]
where any part which is singular in this expansion is called a contraction (we are sort of generalizing the Wick theorem) and they will be represented by the connection among fields.
Thus up to regular terms, we can write
\[
    \langle J(z)J(w) \rangle = \frac{1}{(z-w)^2} = \contraction{}{J(z)}{}{J(w)} J(z) J(w),
\]
so that we can write the operator product expansion of \(J(z)\) with itself as
\[
    J(z)J(w) = \frac{1}{(z-w)^2} + \text{reg.}(z-w)^n,
\]
where the regular terms are less singular than the leading term, and they can be expanded in a power series in \(z-w\). Thus when we are computing the two point function of \(J(z)\) with itself, we are basically picking up the singular part of the operator product expansion, which is the contraction of the two fields. The regular part does not contribute to the two-point function, since it is regular when \(z\) approaches \(w\), and thus it does not have any singularity.\TODO{check.}
[...]
Then if we compute \(T(z)\) we can find
\[
    T(z) = - \frac{1}{2} : (\partial \phi(z))^2 : = \frac{1}{2} : J(z)^2 :,
\]
[...]
so we have to subtract the diverging part, which is basically the normal ordering of the fields (we are basically subtracting the vacuum expectation value of the product of fields, which is the singular part of the operator product expansion)...
    [...]

If we now compute the operator product expansion of \(T(z)\) with the conserved current \(J(w)\), we can find
\[
    T(z)J(w) = \frac{1}{2} : J(z)^2 : J(w),
\]
where if we apply the Wick theorem (contracting in all possible ways the fields) we can find\TODO{add normal ordering ::}
\[
    T(z)J(w) = \frac{1}{2} \contraction{}{J(z)}{}{J(z)} J(z) J(z) J(w) + \frac{1}{2} \contraction{}{J(z)}{}{J(w)} J(z) J(z) J(w),
\]
but since inside the normal ordering (the fields at most commute) we have only regular terms, the first term does not contribute, and we are left with
\[
    T(z)J(z) = \frac{1}{2} \contraction{}{J(z)}{}{J(w)} J(z) J(z) J(w) = \frac{J}{(z-w)^2} + \text{reg.}
\]

[\textit{OPE draw as z to w (points) }]

we know that the ope for two generic operators has to be of the form
\[
    \dots
\]
so we can expand \(J(z)\) in a Taylor series around \(w\) (since it is analytic in \(z\) and \(w\) is in the domain of analyticity)
\[
    J(z) = \dots
\]
and we can find
\[
    T(z)J(w) = \frac{J(w)}{(z-w)^2} + \frac{\partial J(w)}{z - w} + \text{reg.}
\]
which is a very important result, since it shows that \(J(w)\) is a primary field with scaling dimension \(\Delta = 1\), since the leading term in the OPE has a second order pole (and \(h=1\) is the coefficient in front of the leading term).

For the next OPE, we can compute \(T(z)T(w)\), which is the OPE of the stress-energy tensor with itself, and we can find\TODO{fix wick contractions}
\[
    \begin{aligned}
        T(z)T(w) & = \frac{1}{4} : J(z)J(z) : : J(w)J(w) :                                                                                                         \\
                 & = \frac{1}{4} \contraction{}{J(z)}{}{J(z)} J(z) J(z) : J(w)J(w) : + \frac{1}{4} \contraction{}{J(z)}{J(z)}{J(w)} J(z) J(z) : J(w)J(w) : + \dots
    \end{aligned}
\]
giving us terms (1), (2), and (3) which are the contractions of the fields, and we can compute them one by one:
\begin{enumerate}
    \item The first term gives us only regular terms, since the contraction of \(J(z)\) with itself is regular inside the normal ordering, so it does not contribute to the OPE.
    \item The second term gives us a contribution to the OPE, since the contraction of \(J(z)\) with \(J(w)\) gives us a singular term, and thus it contributes as
          \[
              \contraction{}{J(z)}{J(z)}{J(w)} J(z) J(z) : J(w)J(w) : = \frac{1}{(z-w)^2} : J(w)J(w) :,
          \]
          which can be expanded as \(J(z) = J(w) + (z-w) \partial J(w) + \dots\) in order to get
          \[
              \contraction{}{J(z)}{J(z)}{J(w)} J(z) J(z) : J(w)J(w) : = \frac{:J^2(w):}{(z-w)^2} + \frac{:J(w)\partial J(w):}{z-w} + \text{reg.}
          \]
          One can compute \(2 :J\partial J: = \partial :J^2: \) after some algebra, so we can write the contribution of this term as
          \[
              \contraction{}{J(z)}{J(z)}{J(w)} J(z) J(z) : J(w)J(w) : = \frac{2 :J^2(w):}{(z-w)^2} + \frac{\partial :J^2(w):}{z-w} + \text{reg.}
          \]
          but in the first term we can recognize the stress-energy tensor \(T(w) = \frac{1}{2} :J^2(w):\), and in the second term we can recognize the derivative of the stress-energy tensor \(\partial T(w) = \frac{1}{2} \partial :J^2(w):\), so we can write the contribution of this term as
          \[
              \contraction{}{J(z)}{J(z)}{J(w)} J(z) J(z) : J(w)J(w) : = \frac{2 T(w)}{(z-w)^2} + \frac{\partial T(w)}{z-w} + \text{reg.},
          \]
          and if the term (3) does not give us any contribution, we would have found that \(T(z)\) is a primary field with scaling dimension \(\Delta = 2\), but we will see that this is not the case.
    \item [...]

\end{enumerate}
\[
    T(z)T(w) = \frac{1}{2} \frac{1}{\vert z-w \vert^2} + \frac{2 T(w)}{(z-w)^2} + \frac{\partial T(w)}{z-w} + \text{reg.},
\]
so we have an extra term in the OPE, which is not the transformation of a primary field, and thus \(T(z)\) is not a primary field, but it is a secondary field.
In general this ope in general has an arbitrary constant \(c\) in front of the first term, which is called the \textbf{central charge} of the theory, and it is a very important parameter which characterizes the conformal field theory. In our case we have \(c=1\), but in general it can be any real number, and it can be negative as well.

We do not know yet the physical meaning of the central charge, but we will see that it is related to the number of degrees of freedom of the theory, and it is also related to the anomaly of the theory when we try to quantize it. It is caracteristic of the theory, since it comes from the OPE of the stress-energy tensor with itself, which is a fundamental field in the theory (coming directly from the lagrangian of the theory), and thus it is a parameter which characterizes the theory.

The classification of conformal field theories is the idea to find numbers such as the central charge which can determin univocally the theory, and thus to classify all the possible conformal field theories. This is not possible in general, but for example the \textbf{minimal models} (\(0<c<1\)) are a class of conformal field theories which are completely classified by the central charge and the scaling dimensions of the primary fields, and also completely solved, since we can compute all the correlation functions of the theory. The ising model...

\subsection{The Free Fermion in Two Dimensions}